\chapter{Conclusiones}

\begin{itemize}
    \item Pudimos observar la variación de parámetros correspondientes debido a factores que afectan las mediciones, como los procesos de fabricación, irregularidades en las fuentes de tensión y resistencia de los circuitos experimentales.

\item Además, fue indispensable para comprender como el correcto uso e interpretación de los datasheets (hojas de datos) es esencial para seleccionar el diodo adecuado para cada aplicación y evitar fallas por sobrecorriente, sobre temperatura o mal dimensionamiento del circuito. Gracias a estas hojas de datos se puede diseñar de forma segura y eficiente, asegurando que el componente funcione dentro de sus límites eléctricos y térmicos.

\item Como método didáctico, tras la elaboración de este trabajo, pudimos afianzar los conocimientos vistos en el apartado teórico, los cuales fueron vitales al momento de realizar las mediciones en el laboratorio ya que pudimos realizarlas eficientemente sin sacrificar ningún dispositivo y/o instrumental.
\end{itemize}

\section{instrumentos utilizados}
Multímetros empleados para la medición de tensión y corriente:
\begin{itemize}
    \item ProsKit MT-1210
    \item UNI-T UT33A+
    \item UNI-T UT33B+    
\end{itemize}

Instrumental de laboratorio: 
\begin{itemize}
    \item Fuente de tensión DC F-92
\end{itemize}

Transistor bipolar
\begin{itemize}
    \item NPN 548B
\end{itemize}

Datasheet transistor utilizado
\begin{itemize}
  \item https://www.alldatasheet.com/datasheet-pdf/view/11557/ONSEMI/BC548B.html
\end{itemize}

\section{Simulaciones}
\label{con:simulaciones}
Se pueden encontrar todos los archivos de simulación utilizados en este informe en https://github.com/Corte0/dispositivos/tree/main/TP3/informe/simulaciones

\section{Rúbrica}

\begin{table}[H]
  \centering
  \begin{tabular}{|c|c|c|}
    \hline
    \rowcolor[gray]{0.8}
    \textbf{Tarea} & \textbf{Puntuación} & \textbf{Máximo} \\
    \hline
    Presentación del informe & & 30 \% \\
    \hline
    Explicación del proceso de medición & & 40 \% \\
    \hline
    Defensa de las conclusiones & & 30 \% \\
    \hline
    \textbf{Total} & & \textbf{100 \%} \\
    \hline
  \end{tabular}
\end{table}

\section{Planilla de seguimiento}

\begin{table}[H]
  \centering
  \resizebox{\textwidth}{!}{%
  \begin{tabular}{|c|c|c|c|c|c|c|c|}
    \hline
    \multicolumn{1}{|p{2.5cm}|}{\centering \textbf{Versión} \\ \textbf{del TP}} &
    \multicolumn{1}{p{2.2cm}|}{\centering \textbf{Fecha de} \\ \textbf{inicio}} &
    \multicolumn{1}{p{2.5cm}|}{\centering \textbf{Revisado} \\ \textbf{por}} &
    \multicolumn{1}{p{4cm}|}{\centering \textbf{Resumen Observaciones /} \\ \textbf{Correcciones / Propuestas}} &
    \multicolumn{1}{p{4cm}|}{\centering \textbf{Fecha de} \\ \textbf{retroalimentación enviada}} &
    \multicolumn{1}{p{3cm}|}{\centering \textbf{Cambios realizados por} \\ \textbf{JTP (sí/no)}} &
    \multicolumn{1}{p{3cm}|}{\centering \textbf{Fecha de entrega de} \\ \textbf{versión corregida}} &
    \multicolumn{1}{p{3cm}|}{\centering \textbf{Aprobado por jefe de} \\ \textbf{cátedra (sí/no)}} \\
    \hline
    1.0 (2025) & & & & & & & SI \\
    \hline
  \end{tabular}
  }
\end{table}
